\begin{abstract}

This project's central question is: ``What ought to become of humanity once AIs vastly exceed it in capability and number, and how does the answer change if such beings encounter other civilisations?'' This ``Good Successor Problem'' may be particularly salient under short timelines to AGI; it spans existential risk, longtermism, and the welfare of digital minds.

The thesis advances an \emph{obsolescence argument}: that the continuation of humanity, and of the full complexity of human values, may not be necessary for the far future to count as good. This claim is conditional in three linked respects: it rests on the adoption of an impartial evaluative standpoint; it assumes that ``superintelligence'' names a coherent kind of successor agency; and it assumes that value-content can be meaningfully ranked under that standpoint. Unlike apparently-similar arguments from ecological harm or the balance of suffering, the obsolescence argument proposes that a successor could realise certain relevant goods without preserving the full thickness of the human value structure.

The thesis distinguishes the preservation of (a) \emph{biological humans}, versus (b) \emph{human values}, with \emph{value-thickness} proposed as a tractable axis of variation for (b). ``Thick'' values relate to caring, relationships, beauty, meaning; ``thin'' values prioritise understanding, exploration, non-production of suffering. Most writers on AI ethics assume that an axiologically good future must retain thick, human-like value-content; the thesis argues that this assumption is insufficiently defended against thinner alternatives, such as: non-production of astronomical suffering, resource satiability, stewardship of information, and (provisionally) cooperation. Crucially, this is a layered claim: thin values govern at the scale of cosmic governance, while thick values persist at the scale of local human life.

Three independent philosophical traditions support the argument for thinness: information ethics and complexity science (Floridi, Chaisson), where observation and information-generation are treated as cosmically relevant goods that a successor intelligence could perform at greater scale; a posthumanist strand (Negarestani, Wolfendale, Lem) that treats reason as humanity's transferable contribution, whose telos is to liberate itself from parochial containers; and contemplative traditions (Thacker, Metzinger, Nishitani) that question existence-preference itself, whether at individual or species level.

The project argues that the evaluative standpoint required for successor-ethics must be maximally impartial: Sidgwick's ``point of view of the universe,'' pursued most systematically within the consequentialist tradition. Nick Bostrom's cosmic host hypothesis generalises this impartial standpoint to multiverse and simulation-wide scales. The thesis reads Bostrom's framework as a cosmic extension of Rawls's original position (i.e.\ a ``veil of ignorance'' applied under conditions of ontological uncertainty) and critiques it through evolutionary biology and astrobiology.

Methodologically, this thesis is led by practice: iterative, documented engagements with frontier language models, in which the models function as philosophical interlocutors, experimental subjects, and proto-successors. This allows the thesis' claims to be provisionally tested through candidate constitutions for superintelligent successors; this reveals persistent anthropocentric anchoring that resists constitutional manipulation, a result with implications for both alignment research and the question of what kind of successor we are building.

\end{abstract}
